\chapter{Stylingen og den planlagde foreldinga}

\begin{motto}
Skaden var ikkje at nokon misbrukte designet.\\ Skaden var at designet vart utført kompetent — langs den eine aksen som var sal.
\end{motto}

Medan Europa drøymde om at forma skulle følgje funksjonen, og Ulm drøymde om at ho skulle følgje ein metode, fann amerikansk industridesign mellom krigane ut kva forma faktisk følgde i ein marknad: ho følgde salet. Dette er kapitlet der designet held opp å late som, og vert det det heile tida hadde høve til å bli i fråvêret av eit fundament som kunne seie noko anna — eit reiskap for å få folk til å kjøpe. Og det er kapitlet der bokas tese om \emph{skade} får sitt klåraste tidlege prov. For det som skjer her, er ikkje at gode formgjevarar vert lurte til å gjere vondt. Det som skjer er at dyktige formgjevarar gjer presis det dei vert bedne om, kompetent, langs den eine aksen oppdraget spesifiserer; og at aksen er sal, ikkje bruk, ikkje haldbarheit, ikkje verda. Skaden er ikkje misbruk av faget. Skaden er faget, utført godt, utan eit fundament som tvingar fram fleire omsyn enn det eine kunden bad om.

\section{Strømlinja som stil}

Den fyrste generasjonen amerikanske industridesignarar kom ikkje frå handverket eller akademiet, men frå reklamen og teateret, og det merkast på alt dei gjorde. Raymond Loewy, Norman Bel Geddes, Walter Dorwin Teague, Henry Dreyfuss — dei var menn som forstod at form sel, og som selde det. Loewy gjorde det til ein leveregel og ein tittel: aldri la det gode vere godt nok, \emph{Never Leave Well Enough Alone}\autocite{loewy1951}. Bel Geddes, med bakgrunn frå scenografien, måla i \emph{Horizons} eit framtidsbilete av strømlinja tog, skip og fly som rørte ein heil generasjon, sjølv om dei fleste av forslaga hans aldri vart bygde\autocite{geddes1932horizons}. Den nye yrkesgruppa selde ikkje berre produkt; ho selde framtida som eit bilete, og Bel Geddes var dens fremste teiknar. Det er talande at desse mennene kom frå reklame og scene snarare enn frå verkstaden. Dei hadde lært yrket sitt der oppgåva er å vekkje eit ynske, ikkje å løyse eit teknisk problem, og dei tok denne forståinga med seg inn i industridesignet. Forma var for dei eit kommunikasjonsmiddel, eit løfte retta mot kjøparen sitt blikk, og spørsmålet var ikkje fyrst og fremst om tingen verka, men om han selde. I denne yrkeskulturen var det aldri tale om at funksjonen skulle bestemme forma; forma skulle bestemme salet, og det var ein heilt annan og langt meir presis ærend enn den europeiske retorikken om funksjon nokon gong var.

Den synlege signaturen deira var strømlinja. Frå tridvuåra fekk lokomotiv, bilar, kjøleskap, brødristarar og blyantkvessarar dei same mjuke, dropeforma linjene, som om dei skulle skje gjennom luft. Men eit kjøleskap rører seg ikkje, og ein blyantkvessar har inga aerodynamisk last. Strømlinja var ikkje funksjon; ho var \emph{stil}, ein estetikk lånt frå farten og festa på alt som skulle seljast som moderne. Det er verdt å stogge ved kor reint dette dømet er, for det vender den europeiske retorikken på hovudet. Funksjonalismen i Europa hevda at forma følgde funksjonen, og dreiv i røynda med stil. Den amerikanske strømlinja gjorde det same — dreiv med stil — men utan å hevde noko anna. Ho lånte rett nok eit funksjonsspråk, aerodynamikken sitt, men berre som teikn, som ein måte å seie «framtid» og «fart» på ein brødristar som korkje flyg eller fer. Ingen i bransjen trudde for alvor at den dropeforma brødristaren rista brød betre. Dei visste at han \emph{selde} betre, og det var heile poenget. Der europearen narra seg sjølv, narra amerikanaren berre kunden, og var seg det fullt medviten. Jeffrey Meikle har dokumentert denne perioden grundig, og syner kor medvite strømlinja vart nytta som eit kommersielt verkemiddel — ein måte å få fjorårsmodellen til å sjå gamal ut, og årets til å sjå uunnverleg ut\autocite{meikle1979}. Her ser vi den same gliden som i Europa, frå funksjon til smak, men no utan eit fnugg av dårleg samvit, fordi formålet var ope erklært: forma skulle selje. Det amerikanske designet var meir ærleg om sin eigen estetisme enn den europeiske modernismen, nett fordi det aldri hevda å vere noko anna enn handel.

\section{Foreldinga som forretningsmodell}

Så kjem steget som gjer dette til meir enn ei historie om stil. Dersom forma sin oppgåve er å selje, og kunden alt eig produktet, må produktet på eitt eller anna vis gjerast utdatert, slik at kunden kjøper på nytt. Dette er \emph{den planlagde foreldinga}, og ho vart ikkje funnen opp av kyniske unntak, men forfekta opent av framståande menn i bransjen som ein sunn motor for økonomien. Bilindustrien synte vegen med den årlege modellendringa: ikkje ei betre bil kvart år, men ei \emph{annleis} bil, slik at fjorårets, framleis fullt brukande, vart sosialt forelda. Designaren Brooks Stevens gjorde uttrykket til sitt og sa det reint ut — at oppgåva var å plante i kjøparen ynsket om å eige noko litt nyare, litt betre, litt før det var naudsynt.

Legg merke til kva slags innovasjon dette er. Det er ikkje teknisk foreldring, der det nye verkeleg gjer meir enn det gamle. Det er korkje slitasje som tvingar fram utskifting. Det er \emph{stilforeldring}: forma vert endra med vilje, slik at den eldre forma signaliserer «i går», og eigaren kjenner ynsket om å skifte ut noko som enno verkar. Designet er motoren i dette, og berre designet kan vere det, for det er forma — ikkje funksjonen — som ber signalet om gammalt og nytt. Vance Packard skreiv den store samtidskritikken av dette i \emph{The Waste Makers} i 1960, og han skilde skarpt mellom foreldring av funksjon, av kvalitet og av ynske, og synte at det var den siste industrien hadde gjort til kunst\autocite{packard1960}. Giles Slade har seinare ført historia fram til vår eiga tid og synt at den planlagde foreldinga ikkje var ein parentes i mellomkrigsåra, men ein arvestruktur som går rett inn i elektronikken og den digitale tidsalderen — frå nylonstrømpen som vart laga til å rakne, til mobiltelefonen som ikkje let seg reparere\autocite{slade2006}. Meikle har dessutan synt korleis plasten, materialet som meir enn noko gjorde billeg masseproduksjon og utskifting mogleg, sjølv bar ein heil kultur av forgjengelegheit med seg\autocite{meikle2005plastic}.

Det er to slag foreldring som lett vert blanda saman, og skiljet er heile poenget. Den eine er \emph{teknisk} eller fysisk: tingen vert med vilje laga slik at han sviktar etter ei viss tid — pæra som er konstruert til å brenne ut, strømpen som er laga til å rakne. Den andre er \emph{psykologisk} eller stilmessig: tingen verkar like godt som før, men er gjort umoderne, slik at eigaren skammar seg over han og vil ha den nye. Det fyrste slaget krev ingeniørkunst og er ofte ulovleg eller i alle høve lett å fordømme. Det andre slaget krev design, og berre design, for det er forma — ikkje mekanikken — som ber bodskapen om gammalt og nytt. Det er difor den planlagde foreldinga i si reine form er eit \emph{designproblem} og ikkje eit ingeniørproblem: oppgåva er å lage to ting som verkar like godt, men der den eine ser ut som framtida og den andre som fortida. Berre formgjevaren kan løyse den oppgåva, fordi berre forma talar det språket. Faget eig denne skaden på ein måte det ikkje eig nokon annan, for han kan ikkje utførast utan fagkunnskapen sjølv.

\section{Skaden som kompetanse}

Her må eg vere presis, for poenget er lett å misforstå. Det eg påstår er ikkje at desse formgjevarane var dårlege. Tvert om: dei var framifrå. Den planlagde foreldinga er eit prov på \emph{kompetent} design, ikkje udugeleg. Å få ein heil befolkning til å ynskje seg det dei alt har, i ein litt nyare form, og å gjere dette år etter år, krev djup innsikt i form, i psykologi, i kva eit teikn signaliserer. Det er fagkunnskap av høg orden. Skaden ligg ikkje i at fagkunnskapen sviktar; skaden ligg i at fagkunnskapen lukkast — fullkome, langs den eine aksen den vart sett til å tene. Éin eingongskopp krev tolv sekund å produsere, tjue minutt å bruke, og fire hundre år å bryte ned. Det er ikkje dårleg design. Det er presis design; presis for éin akse, og blind for alle dei andre.

Dette er kjernen i bokas tese om skade, og perioden her er staden ho fyrst vert udiskutabel. Eit fag utan eit fundament som tvingar fram fleire omsyn enn det eine oppdraget spesifiserer, vil byggje det skadelege like dugande som det gode, fordi det manglar språket til å skilje dei. Den amerikanske industridesignaren hadde ingen intern standard som kunne seie at det å gjere ein brukande ting forelda er eit overgrep mot noko — mot materialet, mot kjøparen, mot kloden — fordi faget aldri hadde bygt eit omgrep om kva ei form skuldar utover salet. Det fanst ikkje noko i fagkunnskapen sjølv som kunne seie nei. Difor vart svaret ja, om og om att, og kvart ja var fagleg upåklageleg.

Og legg merke til kvar lina peikar. Den eingongskulturen som tek til her, med koppen og strømpa og fjorårsmodellen, er det same mønsteret som seinare skal nå si reine form i ekstraksjonsøkonomien: system bygde for å fange, halde og tøme — merksemd, data, lojalitet — kompetent optimerte langs den eine aksen oppdraget gav. Den planlagde foreldinga er ekstraksjonsøkonomien sin oldefar. Skiljet er berre at den seinare varianten ekstraherer frå mennesket sjølv heller enn frå lommeboka. Strukturen er den same: design utført godt, for éin akse, av folk med fagkunnskap og prisar, utan eit fundament som kunne kalle det noko anna enn vellukka.

Det er verdt å sjå presist korleis arven går vidare, for samanhengen er ikkje berre ein lausleg likskap. Den planlagde foreldinga lærte faget tre ting, og alle tre vart med inn i den digitale tidsalderen. Det fyrste er at oppdragsgjevaren, ikkje brukaren, er den forma tener: kunden til formgjevaren er den som betaler, og brukaren er berre den forma skal verke på. Det andre er at åtferd kan formgjevast — at ein gjennom forma kan styre kva folk ynskjer, kjøper og kastar, ikkje berre kva dei kan. Det tredje er at suksess er målbar i éin storleik, salet, og at alt anna er støy. Set desse tre lærdomane inn i ein skjerm i staden for ein brødristar, og byt salet ut med tid brukt og data hausta, og du har skjermøkonomien sin grunnlov, ferdig skriven alt i mellomkrigsåra. Faget bar med seg, frå strømlinja til strøymetenesta, ein komplett oppskrift på korleis ein formgjev menneskeleg åtferd til oppdragsgjevaren si vinning — og bar samstundes med seg det same hòlet der eit fundament skulle ha vore, det same fråvêret av eit språk som kunne seie at noko av dette var gale.

\section{Dreyfuss og det sjeldne motdømet}

Det ville vere urettferdig å late perioden stå som ei rein anklage, for han rommar òg det sjeldnaste i heile denne boka: eit døme på at eit \emph{målbart} fundament faktisk disiplinerte forma. Henry Dreyfuss skil seg frå dei andre, ikkje fordi han stod utanfor handelen — han teikna for handelen heile livet — men fordi han la til eitt omsyn som let seg måle, og lét det forme arbeidet hans. Det omsynet var menneskekroppen. I \emph{Designing for People} la han fram ein praksis der formgjevaren tok utgangspunkt i den verkelege brukaren, hans kropp, hans grenser, hans rørsler\autocite{dreyfuss1955}. Og i \emph{The Measure of Man} gav han faget det næraste det har kome ein objektiv målestav: antropometriske tabellar og standardiserte menneskefigurar, Joe og Josephine, som la kroppen sine faktiske mål til grunn for kva ei god form var\autocite{dreyfuss1960measure}.

Sjå kva som hender når eit slikt fundament er til stades. Plutseleg kan eit formval vere \emph{feil} på ein måte som ikkje er smak. Eit handtak sit for høgt; ein knapp er for liten for fingeren; ein avstand er større enn rekkjevidda til ein person i femte percentil. Dette er ikkje meiningar. Det er målbare brot mot ein standard utanfor formgjevaren sjølv, og difor kan det rettast, prøvast, etterprøvast. Ergonomien gav faget, på eitt avgrensa felt, nett det heile faget mangla: ein storleik å måle forma mot, som ikkje var den som teikna. Der strømlinja og foreldinga svara på ein akse som tente seljaren, svara ergonomien på ein akse som tente brukaren sin kropp, og den aksen var hard, talfest og felles.

Men legg merke til både kor mektig og kor avgrensa dette motdømet er, for begge sider er lærerike. Mektig: det syner at design \emph{kan} ha eit fundament, at det ikkje er noko prinsipielt umogleg ved å disiplinere forma med ein målbar standard — heile boka sin positive påstand er at dette må gjerast for forma som heile, ikkje berre for kroppen. Avgrensa: ergonomien disiplinerer berre den eine aksen der kroppen møter tingen. Han har ingenting å seie om alt det andre forma òg gjer — om foreldinga, om materialet sin lagnad, om kva tingen gjer med kulturen eller kloden. Dreyfuss kunne gjere eit handtak rett for handa og samstundes gjere produktet forelda til neste år; dei to omsyna låg på heilt ulike akse, og berre det eine hadde ein målestav.

Spør kvifor ergonomien lukkast der resten av faget famla, og svaret peikar rett mot bokas tese. Ergonomien hadde noko ingen annan del av faget hadde: eit \emph{objekt utanfor formgjevaren} som ikkje let seg snakke bort. Menneskekroppen er som han er. Han har mål ein kan ta, grenser ein kan finne, og desse måla rettar seg ikkje etter kva formgjevaren føretrekkjer. Difor kunne det vekse fram ein kumulativ kunnskap — tabellar som vart betre, standardar som vart presisare — for det fanst noko utanfor faget å bryne kunnskapen mot. Alt det andre faget skulle uttale seg om, mangla eit slikt objekt. «God form» har ingen kropp å målast mot; «riktig uttrykk» har inga tabell. Lærdomen er presis: eit fundament oppstår der, og berre der, det finst ein storleik utanfor utøvaren som vala kan prøvast mot. Ergonomien fann ein slik storleik for den eine aksen. Bokas påstand er at ein slik storleik må finnast — eller byggjast — for forma som heile, og at faget aldri gjorde det. Ergonomien er difor ikkje fundamentet faget treng. Han er noko meir lærerikt: eit lite, verkeleg bevis på at eit fundament er mogleg, midt i ein periode som elles syner kor øydeleggjande det er å mangle eitt.

\vignett

Strømlinja og foreldinga selde framtida som ei vare i Amerika. Men det fanst ein annan måte å selje god form på, reinare og meir moralsk i tonen, og han voks fram i nord, kring eit ord som skulle bli eit av dei mektigaste merkevarene faget nokon gong laga: «skandinavisk design». Det lova ikkje utskifting, men varigheit; ikkje sløsing, men måtehald og demokrati. Lat oss sjå kva myten lova, og kva han eigentleg kvilte på.

\vidarelesing{\fullcite{loewy1951}; \fullcite{geddes1932horizons}; \fullcite{meikle1979}; \fullcite{packard1960}; \fullcite{slade2006}; \fullcite{dreyfuss1955}; \fullcite{dreyfuss1960measure}}
